\section{The CPU octree data structure}
\label{sec:octreedata}

\subsection{Node layout}

An \code{OctreeNode} is a compact record:
\begin{center}
\begin{tabular}{@{}p{0.32\linewidth}p{0.60\linewidth}@{}}
\toprule
Field & Purpose\\
\midrule
\code{float sdf[8]} & eight corner samples (may contain $+\infty$)\\
\code{int brushIndex} & material index, or the discard sentinel\\
\code{int textureType} & triplanar / projection / custom mapping\\
\code{glm::vec2 data} & mapping scale or authored UV\\
\code{glm::vec3 hsv} & tint carried to the vertex payload\\
\code{uint8_t lod} & stored ladder level, +1 shifted; $0$ = unset\\
\code{uint8_t chunkLod} & chunk-rung level, +1 shifted; $0$ = not a chunk\\
\code{uint8_t bits} & leaf flag, chunk flag, type\\
\code{uint blockId} & index of the child block in the allocator\\
\code{uint version} & bump on every published revision\\
\bottomrule
\end{tabular}
\end{center}

The +1 shift is a storage trick: a single byte can represent ``unset''
(0) and levels 0--254 without a separate sentinel.  Every propagator uses
the same convention:
\begin{equation}
  lod_{\text{parent}} = \min_i lod_{\text{child},i} + 1,
  \qquad
  chunkLod_{\text{parent}} = \min_i chunkLod_{\text{child},i} + 1 .
\end{equation}
The GPU shader compares the stored level directly against the target
level; no shift is applied in the shader.

Children are stored in a separate \code{ChildBlock}, an eight-entry array
allocated from an \code{OctreeAllocator}.  A leaf has no block; an internal
node has a block with up to eight children.  The allocator is a pool of
blocks and nodes; freeing a subtree returns blocks to the pool and bumps
the generation of the parent.  Because node addresses are recycled, every
publish/delete handler must be paired: a new node at a recycled address
must not inherit the previous occupant's pending mesh.  The engine's
change handler explicitly reports the whole deleted subtree before
\code{clear()} releases the children.

\subsection{Cube corners and edges}

The corner convention is bit-based:
\begin{center}
\begin{tabular}{@{}c|ccc@{}}
\toprule
$i$ & $x$ & $y$ & $z$\\
\midrule
0 & 0 & 0 & 0\\
1 & 0 & 0 & 1\\
2 & 0 & 1 & 0\\
3 & 0 & 1 & 1\\
4 & 1 & 0 & 0\\
5 & 1 & 0 & 1\\
6 & 1 & 1 & 0\\
7 & 1 & 1 & 1\\
\bottomrule
\end{tabular}
\end{center}
The twelve edges are
\begin{equation}
\begin{aligned}
&(0,1),(1,3),(2,3),(0,2),\\
&(4,5),(5,7),(6,7),(4,6),\\
&(0,4),(1,5),(2,6),(3,7).
\end{aligned}
\label{eq:edges}
\end{equation}
This ordering is load-bearing: the QEF loop, the GPU scan loop, the
Blender exporter, and the CPU reference all iterate the same twelve edges
in the same order.  Changing it changes the floating-point accumulation
order and therefore the last bits of a vertex position.

\begin{figure}[htbp]
\centering
\begin{tikzpicture}[scale=0.95]
  % back face
  \draw[grid] (0,0) rectangle (3,3);
  % front face offset
  \draw[grid] (1.1,1.1) rectangle (4.1,4.1);
  \foreach \p in {(0,0),(3,0),(0,3),(3,3)} \draw[grid] \p -- ++(1.1,1.1);
  % corners
  \foreach \x/\y/\lab in {0/0/0,3/0/1,0/3/2,3/3/3,1.1/1.1/4,4.1/1.1/5,1.1/4.1/6,4.1/4.1/7}
    \fill[accent] (\x,\y) circle (2pt) node[below right,font=\scriptsize]{\lab};
  % edges highlighted
  \draw[accent2,thick] (0,0)--(3,0) node[midway,below,font=\scriptsize]{(0,1)};
  \draw[accent2,thick] (3,0)--(4.1,1.1);
  \draw[accent2,thick] (0,0)--(1.1,1.1);
  \draw[accent2,thick] (1.1,1.1)--(4.1,1.1);
  % axes
  \draw[-{Latex[length=2mm]},ink!70] (0,-0.6)--(1,-0.6) node[right,font=\scriptsize]{$x$};
  \draw[-{Latex[length=2mm]},ink!70] (0,-0.6)--(0,0.4) node[above,font=\scriptsize]{$y$};
  \draw[-{Latex[length=2mm]},ink!70] (0,-0.6)--(-0.5,-1.2) node[below,font=\scriptsize]{$z$};
\end{tikzpicture}
\caption{Corner numbering and the first four edges.  The cube is drawn with
$z$ receding; the corner bits are $(x,y,z)$ with bit 2 $=x$, bit 1 $=y$,
bit 0 $=z$.  The GPU table stores corners in two \code{vec4}s and edges as
twelve integer pairs.}
\label{fig:corners}
\end{figure}

\subsection{Types and sentinels}

Each node has a \code{SpaceType}: \code{Empty}, \code{Surface}, or
\code{Solid}.  The classification uses \code{SDF::eval}:
\begin{equation}
  \operatorname{eval}(\mathbf{d}) =
  \begin{cases}
    \text{Empty}, & \text{no finite negative corner and no finite positive corner},\\
    \text{Solid}, & \text{all finite corners negative},\\
    \text{Surface}, & \text{otherwise (both signs present)}.
  \end{cases}
\end{equation}
An all-$+\infty$ cell is Empty.  A cell with only positive finite corners
is Empty.  A cell with only negative finite corners is Solid.  A cell with
both is Surface.  The edge scan is stricter: it tests the \emph{sign pair}
of each edge directly, so a surface cell whose only negative corner is
diagonally opposite a positive one may have no active cube edge; the scan
simply emits nothing.

\subsection{Level of detail and chunk ladder}

Two distinct levels live on a node:
\begin{description}
\item[\code{lod}] the simplification ladder level.  A leaf at the
  tessellation frontier has $lod=1$; a parent whose eight children can be
  collapsed has $lod=\min child + 1$; the root has the largest value.
\item[\code{chunkLod}] the chunk rung.  A node is a chunk root when its
  cube is between $\Delta/2$ and $\Delta$; its immediate descendants are
  marked $chunkLod=1$, their children $2$, and so on.  The renderer's LoD
  gate consumes this value.
\end{description}
The chunk size $\Delta$ is fixed per scene (for example $512$ world
units).  A chunk root at level 0 is the finest published rung; its
ancestors form the coarser rungs.  The renderer keeps exactly one rung per
region visible, which is why the mesher publishes every rung rather than
only the finest one.

\begin{figure}[htbp]
\centering
\begin{tikzpicture}[scale=0.85]
  % tree
  \node[box,minimum width=17mm] (root) at (0,3.2) {root\\$lod=N$};
  \foreach \i/\x in {0/-2.4,1/-0.8,2/0.8,3/2.4} {
    \node[box,minimum width=14mm] (c\i) at (\x,1.7) {chunk\\$lod=N-1$};
    \draw[arrow] (root) -- (c\i);
  }
  \foreach \i/\x in {0/-3.2,1/-2.4,2/-1.6,3/-0.8,4/0.0,5/0.8,6/1.6,7/2.4} {
    \node[box,minimum width=11mm,font=\scriptsize] (g\i) at (\x,0.2) {leaf\\$lod=1$};
  }
  % ladder bar
  \fill[accent!15] (-4.2,-1.0) rectangle (4.2,-0.65);
  \foreach \x/\t in {-3.6/0,-2.4/1,-1.2/2,0/3,1.2/4,2.4/5,3.6/6}
    \node[font=\scriptsize] at (\x,-0.82) {\t};
  \node[font=\scriptsize,left] at (-4.3,-0.82) {rung $r$:};
  \draw[decorate,decoration={brace,mirror,amplitude=4pt},accent2]
    (-3.9,-1.1) -- (3.9,-1.1) node[midway,below=4pt,font=\scriptsize]{published solid chunks (one draw entry per rung)};
\end{tikzpicture}
\caption{The chunk ladder.  Every surface-bearing cell whose size lies in
$(\Delta/2,\Delta]$ is a chunk root at some rung.  Its children are the
next finer rung, and so on.  The GPU gate walks this dyadic lattice from
the root to choose exactly one visible rung per region.}
\label{fig:ladder}
\end{figure}

\subsection{Memory model}

The octree is a sparse pointer structure, not a dense grid.  A node is
small (the eight floats dominate) and a block is eight pointers.  For a
reference terrain with about $10^6$ surface cells the node storage is on
the order of tens of megabytes, which is small compared with the GPU
pools.  The expensive part is not the octree itself but the per-chunk
flattened upload: a chunk's neighbourhood can contain several thousand
nodes, each uploaded as a 128-byte record (eight floats, cube, metadata,
HSV/mapping, children).  \Cref{sec:perf} measures this traffic.

\subsection{Concurrency}

\code{Octree::apply} takes a unique lock on \code{treeMutex}; every
iteration (`iterateMultiThreaded`, the exporter walk, the debug cube walk)
takes a shared lock.  The shape recursion forks at a band of cube sizes and
joins its children before returning, so a tree never exposes a partially
written child block to a reader.  The flattening runs under the shared
lock: it reads the tree structure and copies the neighbourhood into a
private \code{OctreeTopologyInput}.  The GPU pass then works on the copy,
so the tree can be edited again while the GPU is still meshing an older
snapshot.  The version recorded in the chunk manager (not the pointer)
decides whether a rebuild is needed.
